%------------------------------------------------------------------------------%
\documentclass[fleqn,utf8,aspectratio=169,12pt]{beamer}

\usepackage{gaal}

\title{Equação da Reta}

%------------------------------------------------------------------------------%
\begin{document}

\addtitle

%------------------------------------------------------------------------------%
\section{Equação Vetorial da Reta}

%------------------------------------------------------------------------------%
\framedgraphic*{Equação Vetorial da Reta}{F-01-Equacao_da_reta-1}{}
\framedgraphic*{Equação Vetorial da Reta}{F-01-Equacao_da_reta-2}{}
\framedgraphic*{Equação Vetorial da Reta}{F-01-Equacao_da_reta-3}{}
\framedgraphic*{Equação Vetorial da Reta}{F-01-Equacao_da_reta-4}{}
\framedgraphic*{Equação Vetorial da Reta}{F-01-Equacao_da_reta-5}{}
\framedgraphic*{Equação Vetorial da Reta}{F-01-Equacao_da_reta-6}{}
\framedgraphic*{Equação Vetorial da Reta}{F-01-Equacao_da_reta-7}{}
\framedgraphic*{Equação Vetorial da Reta}{F-01-Equacao_da_reta-8}{}
\framedgraphic*{Equação Vetorial da Reta}{F-01-Equacao_da_reta-9}{}
\framedgraphic*{Equação Vetorial da Reta}{F-01-Equacao_da_reta-A}{}

%------------------------------------------------------------------------------%
\begin{frame}{Equação Vetorial da Reta}
  \onslide<+->{Uma reta é determinada por}
  \begin{itemize}[<+->]
      \item um ponto pertencente à reta
      \item um vetor não nulo paralelo à reta
            \onslide<+->{(vetor diretor)}
  \end{itemize}
\end{frame}

%------------------------------------------------------------------------------%
\begin{frame}{Equação Vetorial da Reta}
  Dados um \emph{ponto da reta}, $P$,
  \pause
  e um \emph{vetor diretor}, $\vec{v}$

  \pause
  Para qualquer ponto $R$ da reta

  \pause
  o vetor $\overrightarrow{PR}$ é paralelo a $\vec{v}$

  \pause
  Assim, existe $t\in\R$, tal que
  \[
    \pause
    \overrightarrow{PR} = t\vec{v}
  \]
  \pause
  então
  \[
    R = P + t\vec{v}
  \]
\end{frame}

%------------------------------------------------------------------------------%
\begin{frame}{Equação vetorial da reta}
  Se\quad
  \(
    \vec{p} = \vetortri{a}{b}{c}         \qquad \pause
    \vec{v} = \vetortri{v_1}{v_2}{v_3}   \qquad \pause
    \vec{r} = \vetortri{x}{y}{z}
  \)

  \pause
  A equação vetorial
  \[
    \vec{r} = \vec{p} + t\vec{v}
  \]
  \pause
  é escrita como
  \[
    \pause
    \vetortri{x}{y}{z}
    \pause
    \;=\; \vetortri{a}{b}{c}
    \pause
    \;+\; t\; \vetortri{v_1}{v_2}{v_3}
    \pause
    \qquad t \in \R
  \]

  \pause
  $\vec{r}$ é chamado de \emph{vetor posição} e representa um ponto qualquer
\end{frame}

%------------------------------------------------------------------------------%
\section{Equações Paramétricas}

%------------------------------------------------------------------------------%
\begin{frame}{Equações Paramétricas}
  Escrevendo
 \[
    \vetortri{x}{y}{z}
    \;=\; \vetortri{a}{b}{c}
    \;+\; t\; \vetortri{v_1}{v_2}{v_3}
    \qquad t \in \R
  \]
  \pause
  separadamente para cada coordenada,
  \pause
  temos as \emph{Equações Paramétricas da Reta}
  \[
    \pause
    \begin{cases}
      x = a + t\,v_1 \\
      y = b + t\,v_2 \\
      z = c + t\,v_3
    \end{cases}
  \]
  \pause
  $t \in \R$ é o \emph{parâmetro}
\end{frame}

%------------------------------------------------------------------------------%
%------------------------------------------------------------------------------%
\begin{question}
  Determine a equação vetorial e as equações paramétricas da reta que passa pelo ponto
  \[
    P = (2, -1, 3)
  \]
  e possui vetor diretor
  \[
    \vec{v} = \vetortri{-1}{2}{4}
  \]
\end{question}

%------------------------------------------------------------------------------%
\begin{resolution}{Solução}
  A equação vetorial é obtida diretamente de
  \[
    \vec{r} = P + t\vec{v}
  \]
  \[
    \pause
    \vetortri{x}{y}{z} = \vetortri{2}{-1}{3} + t \vetortri{-1}{2}{4}
  \]
  \[
    \pause
    \begin{cases}
      x =  2 -  t \\
      y = -1 + 2t \\
      z =  3 + 4t
    \end{cases}
  \]
\end{resolution}

%------------------------------------------------------------------------------%
\endinput

%------------------------------------------------------------------------------%
\begin{question}
  Determine as equações paramétricas da reta que passa pelos pontos
  \[
    A = (1, 2, -1)
  \]
  \[
    B = (4, -2, 5)
  \]
\end{question}

%------------------------------------------------------------------------------%
\begin{resolution}{Solução}
  Vetor diretor
  \pause
  \[
    \vec{v}
    \pause
    = \overrightarrow{AB}
    \pause
    = B - A
    \pause
    = \vetortri{4}{-2}{5} - \vetortri{1}{2}{-1}
    \pause
    = \vetortri{3}{-4}{6}
  \]

  \pause
  Equação da reta
  \(
    \pause
    \vec{r} = A + t\vec{v}
  \)
  \[
    \pause
    \vetortri{x}{y}{z} = \vetortri{1}{2}{-1} + t\vetortri{3}{-4}{6}
  \qquad\qquad\qquad
    \pause
    \begin{cases}
      x =  1 + 3t \\
      y =  2 - 4t \\
      z = -1 + 6t
    \end{cases}
  \]
\end{resolution}

%------------------------------------------------------------------------------%
\endinput

%------------------------------------------------------------------------------%
\begin{question}
  Determine se o ponto
  \(
    P = (0, 7, 8)
  \)
  pertence à reta
  \[
    r\colon
    \begin{cases}
      x = 2 -  t \\
      y = 1 + 3t \\
      z = 4 + 2t
    \end{cases}
  \]
\end{question}

%------------------------------------------------------------------------------%
\begin{resolution}{Solução}
  \onslide<+->{Para que $P$ pertença à reta,}
  \onslide<+->{deve existir um mesmo valor de $t\in\R$} \\
  \onslide<+->{que satisfaça as três equações}
\begin{columns}
\column{0.3\textheight}
  \begin{align*}
    \onslide<+->{x & = 2 - t \\}
    \onslide<+->{0 & = 2 - t \\}
    \onslide<+->{t & = 2}
  \end{align*}
\column{0.3\textheight}
  \begin{align*}
    \onslide<+->{y & = 1 + 3t \\}
    \onslide<+->{  & = 1 + 3\times 2 \\}
    \onslide<+->{  & = 7}
  \end{align*}
\column{0.3\textheight}
  \begin{align*}
    \onslide<+->{z & = 4 + 2t \\}
    \onslide<+->{  & = 4 + 2\times 2 \\}
    \onslide<+->{  & = 8}
  \end{align*}
\end{columns}

  \onslide<+->{As três coordenadas são satisfeitas pelo mesmo parâmetro,} \\
  \onslide<+->{portanto, $P\in r$}
\end{resolution}

%------------------------------------------------------------------------------%
\endinput


%------------------------------------------------------------------------------%
\section{Equações Simétricas}

%------------------------------------------------------------------------------%
\begin{frame}{Equações Simétricas}
  Se $v_1$, $v_2$ e $v_3$ forem diferentes de zero,
  \pause
  isolamos o $t$ em
  \[
    \pause
    \begin{cases}
      x = a + tv_1 \\
      y = b + tv_2 \\
      z = c + tv_3
    \end{cases}
  \]
  \pause
  obtendo
  \[
    \pause
    t = \frac{x-a}{v_1} \qquad \pause
    t = \frac{y-b}{v_2} \qquad \pause
    t = \frac{z-c}{v_3}
  \]
  \pause
  Para qualquer ponto na reta
  \[
    \pause
    \frac{x-a}{v_1} \;=\;
    \frac{y-b}{v_2} \;=\;
    \frac{z-c}{v_3}
  \]
\end{frame}

%------------------------------------------------------------------------------%
%------------------------------------------------------------------------------%
\begin{question}
  Determine a equação simétrica da reta
  \[
    r\colon
    \begin{cases}
      x =  2 + 3t \\
      y = -1 - 2t \\
      z =  4 +  t
    \end{cases}
  \]
\end{question}

%------------------------------------------------------------------------------%
\begin{resolution}{Solução}
  Isolando o parâmetro em cada equação
  \[
    \pause
    t =  \frac{x-2}{3} \qquad\pause
    t = -\frac{y+1}{2} \qquad\pause
    t = z-4
  \]
  \pause
  temos
  \[
    \frac{x-2}{3}
    = \frac{y+1}{-2}
    = z-4
  \]
\end{resolution}

%------------------------------------------------------------------------------%
\endinput
%------------------------------------------------------------------------------%

%------------------------------------------------------------------------------%
\begin{question}
  Determine uma equação vetorial para a reta
  \[
    \frac{x-1}{2}
    = \frac{y+3}{-1}
    = \frac{z-2}{4}
  \]
\end{question}

%------------------------------------------------------------------------------%
\begin{resolution}{Solução}
  Introduzimos um parâmetro comum
  \pause
  \[
    \frac{x-1}{2}
    = \frac{y+3}{-1}
    = \frac{z-2}{4} \pause
    \;=\; t
  \]
  \pause
  Assim,
  \[
    x =  1 + 2t \qquad\pause
    y = -3 -  t \qquad\pause
    z =  2 + 4t
  \]
  \pause
  Portanto
  \[
    (x, y, z) = (1, -3, 2) +t(2, -1, 4)
  \]

  \pause
  O ponto
  \(
    P = (1, -3, 2)
  \)
  pertence à reta

  \pause
  e
  \(
    \vec{v} = (2, -1, 4)
  \)
  é um vetor diretor
\end{resolution}

%------------------------------------------------------------------------------%
\endinput


%%------------------------------------------------------------------------------%
%\section{Lista Mínima}
%
%%------------------------------------------------------------------------------%
%\begin{frame}{Lista Mínima}
%  \begin{enumerate}
%    \item
%  \end{enumerate}
%
%  \vfill
%  Atenção: A prova é baseada no livro, não nas apresentações
%\end{frame}

%------------------------------------------------------------------------------%
\end{document}
%------------------------------------------------------------------------------%
